% GENERATED FILE. Edit canonical lessons, then run npm run generate:books.
% canonical-source-hash: fnv1a64:266c600dd848ad1c
% canonical-lessons: BN-C518-ludo, BN-C518-pasha, BN-C518-sapludo, BN-C518-shobdojobdo, BN-C518-kerom

\chapter{Games Indoors}
\label{ch:bn-games-indoors}
\hlchaptermodality{518}

\begin{chapteropening}
\textbf{By the end of this chapter:} \emph{I can talk about board games and crosswords.}

Complete the last lesson of chapter 518: I can talk about board games and crosswords.
\end{chapteropening}

\section[\texorpdfstring{\bn{লুডো} (luḍo)}{লুডো (luḍo)}]{\bn{লুডো} (luḍo) — ludo}

\label{lesson:BN-C518-ludo}



\begin{quote}
Before the new one: say the Bengali for hide-and-seek, then the Bengali for a marble.
\end{quote}

\subsection*{You'll want to know: \bn{লুডো}}
\textbf{\bn{লুডো}} — \emph{luḍo} — \textquotedblleft{}ludo\textquotedblright{}.

\begin{quote}
\bn{বৃষ্টির} \bn{দিনে} \bn{আমরা} \bn{লুডো} \bn{খেলি।}
\end{quote}

\begin{quote}
\emph{Brishṭir dine āmrā luḍo kheli.}
\end{quote}

\begin{quote}
On rainy days we play ludo.
\end{quote}

\begin{cousinweb}
\textbf{\bn{লুডো}} (\emph{luḍo}) is borrowed from English \emph{ludo}, the name of the game.
\end{cousinweb}

\begin{grammarlens}[title={What you've built}]
One more word for describing what you lived and felt.
\end{grammarlens}

\subsection*{Your turn}
\begin{itemize}
\raggedright
  \item \emph{Say it:} \emph{luḍo}
  \item \emph{Say it:} \emph{luḍo}, once more
  \item \emph{Say it:} \emph{Brishṭir dine āmrā luḍo kheli.} — On rainy days we play ludo.
  \item \emph{Recall:} say the Bengali for hide-and-seek, then the Bengali for a marble, then say \emph{luḍo} again
\end{itemize}

\begin{tcolorbox}[breakable,colback=teal!4,colframe=teal!35!black,title={Before you move on}]
What does \bn{লুডো} mean? (\textquotedblleft{}Ludo\textquotedblright{}.) Say it once more.
\end{tcolorbox}
\section[\texorpdfstring{\bn{পাশা} (pāshā)}{পাশা (pāshā)}]{\bn{পাশা} (pāshā) — dice; a dice game}

\label{lesson:BN-C518-pasha}



\begin{quote}
Before the new one: say the Bengali for a marble, then the Bengali for ludo.
\end{quote}

\subsection*{You'll want to know: \bn{পাশা}}
\textbf{\bn{পাশা}} — \emph{pāshā} — \textquotedblleft{}dice; a dice game\textquotedblright{}.

\begin{quote}
\bn{গল্পের} \bn{রাজা} \bn{পাশা} \bn{খেলায়} \bn{প্রাসাদ} \bn{হারালেন।}
\end{quote}

\begin{quote}
\emph{Gôlper rājā pāshā khelāy prāsād hārālen.}
\end{quote}

\begin{quote}
In the story, the king lost his palace at dice.
\end{quote}

\begin{cousinweb}
\textbf{\bn{পাশা}} (\emph{pāshā}) is a tadbhava: CDIAL lists Bengali \emph{pāsā}, \textquotedblleft{}a die\textquotedblright{}, under Sanskrit \emph{pāśa} (CDIAL 8132).
\end{cousinweb}

\begin{grammarlens}[title={What you've built}]
One more word for describing what you lived and felt.
\end{grammarlens}

\subsection*{Your turn}
\begin{itemize}
\raggedright
  \item \emph{Say it:} \emph{pāshā}
  \item \emph{Say it:} \emph{pāshā}, once more
  \item \emph{Say it:} \emph{Gôlper rājā pāshā khelāy prāsād hārālen.} — In the story, the king lost his palace at dice.
  \item \emph{Recall:} say the Bengali for a marble, then the Bengali for ludo, then say \emph{pāshā} again
\end{itemize}

\begin{tcolorbox}[breakable,colback=teal!4,colframe=teal!35!black,title={Before you move on}]
What does \bn{পাশা} mean? (\textquotedblleft{}Dice; a dice game\textquotedblright{}.) Say it once more.
\end{tcolorbox}
\section[\texorpdfstring{\bn{সাপলুডো} (sāpluḍo)}{সাপলুডো (sāpluḍo)}]{\bn{সাপলুডো} (sāpluḍo) — snakes and ladders}

\label{lesson:BN-C518-sapludo}



\begin{quote}
Before the new one: say the Bengali for ludo, then the Bengali for dice.
\end{quote}

\subsection*{You'll want to know: \bn{সাপলুডো}}
\textbf{\bn{সাপলুডো}} — \emph{sāpluḍo} — \textquotedblleft{}snakes and ladders\textquotedblright{}.

\begin{quote}
\bn{ছুটির} \bn{দিনে} \bn{আমরা} \bn{সাপলুডো} \bn{খেলি।}
\end{quote}

\begin{quote}
\emph{Chhuṭir dine āmrā sāpluḍo kheli.}
\end{quote}

\begin{quote}
On holidays we play snakes and ladders.
\end{quote}

\begin{cousinweb}
\textbf{\bn{সাপলুডো}} (\emph{sāpluḍo}) is built from \bn{সাপ}, \textquotedblleft{}a snake\textquotedblright{}, and \bn{লুডো}, \textquotedblleft{}ludo\textquotedblright{}: the board game of snakes and ladders.
\end{cousinweb}

\begin{grammarlens}[title={What you've built}]
One more word for describing what you lived and felt.
\end{grammarlens}

\subsection*{Your turn}
\begin{itemize}
\raggedright
  \item \emph{Say it:} \emph{sāpluḍo}
  \item \emph{Say it:} \emph{sāpluḍo}, once more
  \item \emph{Say it:} \emph{Chhuṭir dine āmrā sāpluḍo kheli.} — On holidays we play snakes and ladders.
  \item \emph{Recall:} say the Bengali for ludo, then the Bengali for dice, then say \emph{sāpluḍo} again
\end{itemize}

\begin{tcolorbox}[breakable,colback=teal!4,colframe=teal!35!black,title={Before you move on}]
What does \bn{সাপলুডো} mean? (\textquotedblleft{}Snakes and ladders\textquotedblright{}.) Say it once more.
\end{tcolorbox}
\section[\texorpdfstring{\bn{শব্দজব্দ} (shôbdôjôbdô)}{শব্দজব্দ (shôbdôjôbdô)}]{\bn{শব্দজব্দ} (shôbdôjôbdô) — a crossword}

\label{lesson:BN-C518-shobdojobdo}



\begin{quote}
Before the new one: say the Bengali for dice, then the Bengali for snakes and ladders.
\end{quote}

\subsection*{You'll want to know: \bn{শব্দজব্দ}}
\textbf{\bn{শব্দজব্দ}} — \emph{shôbdôjôbdô} — \textquotedblleft{}a crossword\textquotedblright{}.

\begin{quote}
\bn{বাবা} \bn{রোজ} \bn{সকালে} \bn{পত্রিকার} \bn{শব্দজব্দ} \bn{করেন।}
\end{quote}

\begin{quote}
\emph{Bābā roj sôkāle pôtrikār shôbdôjôbdô kôren.}
\end{quote}

\begin{quote}
My father does the magazine's crossword every morning.
\end{quote}

\begin{cousinweb}
\textbf{\bn{শব্দজব্দ}} (\emph{shôbdôjôbdô}) is a playful modern compound: \bn{শব্দ}, \textquotedblleft{}a word\textquotedblright{}, and \bn{জব্দ}, \textquotedblleft{}caught out, beaten\textquotedblright{}, which is usually traced to Arabic \emph{ḍabṭ}, \textquotedblleft{}seizing\textquotedblright{}: the words that catch you out.
\end{cousinweb}

\begin{grammarlens}[title={What you've built}]
One more word for describing what you lived and felt.
\end{grammarlens}

\subsection*{Your turn}
\begin{itemize}
\raggedright
  \item \emph{Say it:} \emph{shôbdôjôbdô}
  \item \emph{Say it:} \emph{shôbdôjôbdô}, once more
  \item \emph{Say it:} \emph{Bābā roj sôkāle pôtrikār shôbdôjôbdô kôren.} — My father does the magazine's crossword every morning.
  \item \emph{Recall:} say the Bengali for dice, then the Bengali for snakes and ladders, then say \emph{shôbdôjôbdô} again
\end{itemize}

\begin{tcolorbox}[breakable,colback=teal!4,colframe=teal!35!black,title={Before you move on}]
What does \bn{শব্দজব্দ} mean? (\textquotedblleft{}A crossword\textquotedblright{}.) Say it once more.
\end{tcolorbox}
\section[\texorpdfstring{\bn{ক্যারম} (kærôm)}{ক্যারম (kærôm)}]{\bn{ক্যারম} (kærôm) — carrom (a board game)}

\label{lesson:BN-C518-kerom}



\begin{quote}
Before the new one: say the Bengali for snakes and ladders, then the Bengali for a crossword.
\end{quote}

\subsection*{You'll want to know: \bn{ক্যারম}}
\textbf{\bn{ক্যারম}} — \emph{kærôm} — \textquotedblleft{}carrom (a board game)\textquotedblright{}.

\begin{quote}
\bn{দাদু} \bn{আর} \bn{আমি} \bn{সন্ধ্যায়} \bn{ক্যারম} \bn{খেলি।}
\end{quote}

\begin{quote}
\emph{Dādu ār āmi sôndhyāy kærôm kheli.}
\end{quote}

\begin{quote}
My grandfather and I play carrom in the evening.
\end{quote}

\begin{cousinweb}
\textbf{\bn{ক্যারম}} (\emph{kærôm}) is the English name of the game, \emph{carrom}; its earlier origin is not settled.
\end{cousinweb}

\begin{grammarlens}[title={What you've built}]
One more word for describing what you lived and felt.
\end{grammarlens}

\subsection*{Your turn}
\begin{itemize}
\raggedright
  \item \emph{Say it:} \emph{kærôm}
  \item \emph{Say it:} \emph{kærôm}, once more
  \item \emph{Say it:} \emph{Dādu ār āmi sôndhyāy kærôm kheli.} — My grandfather and I play carrom in the evening.
  \item \emph{Recall:} say the Bengali for snakes and ladders, then the Bengali for a crossword, then say \emph{kærôm} again
\end{itemize}

\begin{tcolorbox}[breakable,colback=teal!4,colframe=teal!35!black,title={Before you move on}]
What does \bn{ক্যারম} mean? (\textquotedblleft{}Carrom (a board game)\textquotedblright{}.) Say it once more.
\end{tcolorbox}
